\documentclass[11pt]{article}
\usepackage[margin=1in]{geometry}
\usepackage[T1]{fontenc}
\usepackage{lmodern,amsmath,amssymb,microtype}
\usepackage[colorlinks=true,linkcolor=blue!50!black,urlcolor=blue!50!black]{hyperref}
\usepackage{xcolor}
\hypersetup{pdftitle={A wild cubic coefficient field for torsion specialization}}
\title{A wild cubic coefficient field for torsion specialization}
\author{}\date{}
\begin{document}
\maketitle
\noindent The ProveIt archive
\href{https://github.com/VladimirReshetnikov/ProveIt/blob/5ed40f562710fcc2770bbb463dba528a6727f7ae/docs/incoming/exponent_seven_and_grh_transfers.zip}{\emph{Exponent seven and GRH transfers}}
determines shell obstruction costs over all finite abelian extensions of
$\mathbb Q_p$. It includes all seven dyadic quadratic fields.
We treat the non-Galois cubic $K=\mathbb Q_3(\pi)$, $\pi^3=3$. The
underlying specialization question comes from
\href{https://github.com/lsernau/grh-lean/blob/72fead102f663d5c7a407964a13dc0cf9d977207/grh/src/New/Congruences/Selmer/RootsOfUnity.lean#L300}{the one-variable roots-of-unity lemma}
in the Lean repository. Normalize $v(3)=1$.
\paragraph{The question.} For $Z=\{\zeta-1:\mathrm{ord}(\zeta)=9\}$,
a list of integral polynomials $F_i\in\mathcal O_K[T]$ has witnesses $x_i$
of valuation at least $\rho\ge1/6$, with finite $t_i=v(F_i(x_i))$.
The cost is $\sum_i t_i$. For prescribed $\epsilon\ge0$ we ask for the
least cost forcing $\sum_i\max(0,v(F_i(z))-t_i)>\epsilon$ at every $z\in Z$.

\paragraph{Answer.} The minimum is $1/3$ for $0\le\epsilon<2/9$,
$2/3$ for $2/9\le\epsilon<4/9$, and $1$ for $\epsilon\ge4/9$.
All three can be attained with one polynomial and witness zero.

\paragraph{Reduction.} The one-shell monic-polynomial test in the earlier
ProveIt reports gives the price as $m/6$, where $m$ is the smallest degree
of a monic $P\in\mathcal O_K[T]$ with $v(P(z))>m/6+\epsilon$ for all $z$.
For completeness, for each $F=c\prod(T-\alpha)$ its Gauss valuation is
$G_r(F)=v(c)+\sum_\alpha\min(r,v(\alpha))$.
Only roots with $v(\alpha)=1/6$ can cause positive excess over $G_{1/6}$
at $z$. Each baseline is at least $G_\rho\ge G_{1/6}$. Collect the matching
roots of the list in a monic $P$ of degree $m$. They form a Galois-stable
integral multiset, and contribute $m/6$ to the sum of
$G_\rho(F_i)-G_0(F_i)$. This gives $m/6\le\sum_i t_i$ and bounds total
positive excess by $v(P(z))-m/6$. Conversely a first successful monic
polynomial has all roots matching, because each other root contributes
at most $1/6$ to its value at $z$ and can be removed with its conjugates.
Its witness zero realizes its degree cost.

\paragraph{Proof of the degree test.} Fix $\alpha=\zeta_9-1$ and
$B(T)=T^2+\pi T+\pi$. The conjugates over $K$ of $\alpha$ are all six
points of $Z$. Indeed $K$ has prime degree and is non-Galois
(it does not contain $\zeta_3$), whereas $\mathbb Q_3(\zeta_9)$ is
abelian. Their intersection is $\mathbb Q_3$ and the latter extension
is Galois. Their compositum therefore has degree $18$.
The polynomial $\Phi_9(1+T)$ is Eisenstein, so $v(\alpha)=1/6$.
An exact resultant determines the second basis value,
\[
 \operatorname{Res}_{Y}\bigl(Y^3-3,\,
 \operatorname{Res}_{T}(\Phi_9(1+T),T^2+YT+Y)\bigr)=59049=3^{10}.
\]
All eighteen embeddings preserve the $3$-adic valuation, whence
$v(B(\alpha))=5/9$. The accompanying check computes this norm by two
integer multiplication determinants.

The monic polynomials $T^jB(T)^k$ for $j=0,1$ and $k=0,1,2$ have
consecutive degrees $2k+j=0,\ldots,5$. Their valuation numerators at
$\alpha$, in units of $1/18$ and in increasing degree, are
$0,3,10,13,20,23$. These represent all residue classes modulo six.
Since $v(K^\times)=\frac13\mathbb Z$, the nonzero terms of a linear
combination over $K$ have distinct valuations. A monic polynomial of
degree $2k+j$ must have coefficient one at $T^jB^k$ in this basis,
and therefore has valuation at most $j/6+5k/9$. The monic basis
polynomial attains equality at every conjugate. Subtracting
$(2k+j)/6$ gives the exact excess $2k/9$. Degree six can give infinite
valuation by the shifted cyclotomic polynomial. This yields the three
claimed ranges, including their strict endpoints.
\section*{All primitive $3$-power orders}
The ninth-root calculation extends to every conductor and exhibits a
difference from the eventual stabilization of the entire tolerance
profile over abelian coefficient fields reported in \emph{Exponent seven
and GRH transfers}. Write
\[
 N=2\cdot3^{h-1},\quad r=1/N,\qquad
 \delta_h=\frac12-\frac{5}{3N}\qquad(h\ge2).
\]
On the shell of primitive $3^h$th roots shifted by one, use witness
radius $\rho\ge r$. The optimal blocking price is
\[
 \begin{cases}
 1/3 & 0\le\epsilon<\delta_h,\\
 2/3 & \delta_h\le\epsilon<2\delta_h,\\
 1   & 2\delta_h\le\epsilon .
 \end{cases}                                                     \tag{1}
\]
At $h=1$ the cost is $1$ for all tolerances. The two shifted primitive
cube roots have mutual distance of valuation $1/2$. A monic linear
polynomial cannot have valuation strictly greater than $1/2$ at both,
by the ultrametric inequality applied to their difference. The shifted
cyclotomic polynomial of degree two attains cost $1$.

The two thresholds increase strictly towards $1/2$ and $1$.
In particular the functions of $\epsilon$ differ at every successive
level, although at each fixed $\epsilon$ the sequence of costs eventually
stabilizes. The proof uses standard different and ramification formulas
from J.-P.~Serre, \emph{Local Fields}, Chapters III and IV, including
the compatibility of upper numbering with quotients and lower numbering
with subgroups.

\paragraph{Ramification in the compositum.}
Let $C_j=\mathbb Q_3(\zeta_{3^j})$, $E_j=KC_j$, and $E=E_h$.
As $K$ has degree three and is non-Galois, $[E:C_h]=3$ and
$\operatorname{Gal}(E/K)$ acts transitively on all $N$ test points.
Over $C_1$, the extension $E_1$ is cyclic of degree three, since $C_1$
contains all cube roots of unity. The extension $E_1/\mathbb Q_3$ is
totally ramified of degree six: its ramification index is divisible by
the coprime indices $2$ and $3$. A uniformizer of $E_1$ is
$u=(\zeta_3-1)/\pi$. A nonidentity automorphism over $C_1$ multiplies
$u$ by a primitive cube root of unity. Thus the valuation of $gu-u$
in uniformizer units is $1+3=4$. Both the upper and lower break
of this cubic is $3$.

The lower breaks of $C_h/C_1$ are $3^k-1$ for $1\le k\le h-1$.
Indeed if $v_3(b-1)=k$, the action $\zeta\mapsto\zeta^b$ changes
the uniformizer $\zeta-1$ by an element of normalized valuation $3^k$.
The groups drop in index by three at each break. Integrating the
reciprocal indices gives upper breaks $2,4,\ldots,2(h-1)$.
The Galois group of $E/C_1$ is the product of the two factors just
considered: $E_1\cap C_h=C_1$ since otherwise $K\subset C_h$.
Its upper filtration is the product of the filtrations on the factors.
For detail, an inflated nontrivial character from the cubic factor
has break $3$, and the other nontrivial factor characters have even
breaks. The product of two characters with distinct breaks has break
equal to the larger one, since the smaller character is trivial
between the two breaks. This determines the annihilator of each upper
group, hence the group itself. In particular there is no nontrivial
unramified character, so $E/C_1$ is totally ramified. The upper-to-lower
map takes $3$ to $2+3(3-2)=5$, because the group has index one
through upper level $2$ and index three thereafter up to $3$.
Intersect the lower filtration with $\operatorname{Gal}(E/C_h)$.
This cyclic cubic has just the lower break $5$, and Hilbert's
different formula gives its different exponent $(5+1)(3-1)=12$.

\paragraph{Length of the power-basis order.}
Take $\alpha=\zeta_{3^h}-1$, and write $d_C$ for the
different exponent of $C_h/\mathbb Q_3$. The shifted cyclotomic
polynomial is Eisenstein, so the discriminant of the absolute power
basis has $3$-adic exponent $d_C$. Since its degree stays $N$ over
$K$, the order $\mathcal O_K[\alpha]$ has discriminant exponent
$3d_C$ in $K$. On the other hand the different exponent of
$K/\mathbb Q_3$ is $5$, from the derivative $3\pi^2$. Transitivity
of the different in the two towers through $C_h$ and $K$ gives the
relative different exponent $12+3d_C-5N$ for $E/K$. This extension
is totally ramified, so its relative discriminant has the same
exponent in $K$. Thus the length over $\mathcal O_K$ of the quotient is
\[
 I=\operatorname{length}_{\mathcal O_K}(\mathcal O_E/\mathcal O_K[\alpha])
   =\frac{5N-12}{2}.                                             \tag{2}
\]

\paragraph{The best monic approximation.}
Put $n=N/3$. Among monic integral polynomials of degree $n$, take
$B$ for which $v(B(\alpha))$ is largest. Existence follows from compactness
of $\mathcal O_K^n$ and continuity: the image is compact and cannot contain
zero because $\alpha$ has degree $N$. Put $b=v(B(\alpha))\ge1/3$, using
the candidate $T^n$. We claim $b\notin N^{-1}\mathbb Z$. If it belonged
to this subgroup, write it as $j/N+t/3$ with $0\le j<n$, $t\ge0$.
A suitable unit multiple of $\pi^t\alpha^j$ then cancels the initial
residue of $B(\alpha)$ since both residue fields are $\mathbb F_3$.
This improves the monic approximation, a contradiction.

The $N$ elements $w_{jk}=\alpha^jB(\alpha)^k$,
$0\le j<n$, $0\le k<3$, form an $\mathcal O_K$-basis of its power-basis
order by triangularity in the degrees $m=kn+j$. Their valuations
are pairwise distinct modulo $\tfrac13\mathbb Z$: $b$ has order
three modulo $N^{-1}\mathbb Z$ in the value group $(3N)^{-1}\mathbb Z$,
and the $j/N$ for $j<n$ are distinct modulo $\tfrac13\mathbb Z$.
Every nonzero linear combination over $K$ therefore has valuation
equal to its least term valuation. In particular
$w_{jk}/\pi^{\lfloor3v(w_{jk})\rfloor}$ form an integral basis
of $E/K$. Indeed their scaled valuations $3v$ lie in $[0,1)$ in distinct
classes, so a vector in this basis is integral exactly when all its
coefficients are. It follows that
$I=\sum_{j,k}\lfloor3v(w_{jk})\rfloor$.

Write $\delta=b-1/3$. At degree $m$ the basis valuation is $m/N+k\delta$.
The fractional parts of all $3v(w_{jk})$ are precisely
$0,1/N,\ldots,(N-1)/N$, so their sum is $(N-1)/2$.
Using $\sum_{j,k} k=N$ gives
\[
 I=\sum_{m=0}^{N-1}3m/N+3N\delta-(N-1)/2
   =N-1+3N\delta.
\]
Comparison with (2) yields $\delta=\delta_h$.
A monic polynomial of degree $kn+j<N$ has coefficient one on its leading
basis element and cannot improve its valuation by lower terms.
Thus the maximal excess in that degree is exactly $k\delta_h$,
attained simultaneously on the entire conjugacy orbit by $T^jB^k$.
The single-shell reduction above, with $6$ replaced by $N$, gives (1).
Degree $N$ with infinite valuation is always available.

The independent calculation in \texttt{check\_wild\_tower.py} explicitly
constructs maximizing degree-$N/3$ polynomials for $2\le h\le6$ by
successive residue cancellations. It evaluates each candidate with
the cubic norm and exact reduction modulo the shifted cyclotomic
polynomial. The resulting maximum valuations in units of $1/(3N)$
are $10,40,130,400,1210$. The check uses no ramification-break or
discriminant formula.

\paragraph{Relation to the literature.}
The fields $\mathbb Q_p(\zeta_{p^h},p^{1/p})$ are classical layers of the
false Tate curve extension. Viviani computed ramification and discriminants
in radical cyclotomic extensions. Wang and Yuan constructed explicit
uniformizers in these layers. The maximizing monic polynomials used above
are related to the Okutsu frames underlying integral-basis algorithms.
The application here is the strict positive-loss threshold for a list of
polynomials with separately chosen witnesses.
See \href{https://doi.org/10.5802/jtnb.470}{F.~Viviani, J. Th\'eor.
Nombres Bordeaux 16 (2004)},
\href{https://arxiv.org/abs/2301.09135}{S.~Wang and Y.~Yuan,
\emph{Uniformizer of the False Tate Curve Extension of $\mathbb Q_p$ (II)}},
and \href{https://doi.org/10.1016/j.jnt.2014.07.027}{J.~Guàrdia,
J.~Montes and E.~Nart, \emph{Higher Newton polygons and integral bases}}.
\end{document}

